\section{Introduction}
\label{sec:intro}

Most people who train neural networks never see the machinery that makes it possible: the graph
that records every operation, the chain rule applied backwards through it, the update rule that
turns a gradient into a step. visualML is that machinery written out, small enough to be read in an
afternoon, and complete enough to reach the accuracy one expects on MNIST
(Section~\ref{sec:experiments}).

The project has two halves that share one codebase.
\begin{itemize}
  \item \textbf{The engine} (\code{src/visualml/}): a \code{Tensor} type with reverse-mode automatic
        differentiation (Section~\ref{sec:autodiff}), a set of differentiable operations
        (Section~\ref{sec:ops}), modules and losses (Section~\ref{sec:losses}), and six optimizers
        with a line search (Section~\ref{sec:optim}). It depends on nothing but NumPy.
  \item \textbf{The website} (\code{demos/}), which loads the same Python package into the browser
        with Pyodide (Section~\ref{sec:browser}) and uses it to train networks live, and to take a
        trained one apart number by number.
\end{itemize}

Two rules shape the code. Every differentiable operation is checked against a numerical gradient
before it is used (Section~\ref{sec:verification}). And nothing the website shows is reimplemented
in JavaScript: the page asks the engine, and draws what it answers.

The mathematics in this document is what the author learnt in the first two years of the
Bachelor's Degree in Data Science and Engineering (GCED) at the Universitat Polit\`ecnica de
Catalunya: linear algebra~\cite{upc-alg}, mathematical optimisation~\cite{upc-om}, machine
learning~\cite{upc-aa1} and audiovisual processing~\cite{upc-ipa}. The original papers and books
are cited where each method appears.

% TODO(author): one paragraph on motivation and on what you learnt building it.
